\begin{figure}[H]
\centering
\resizebox{0.94\textwidth}{!}{%
\begin{tikzpicture}[
 stage/.style={draw,rounded corners=3pt,minimum width=33mm,minimum height=21mm,align=center,font=\small,inner xsep=4pt},
 flow/.style={-{Stealth[length=2mm]},very thick,prblue},
 tag/.style={font=\scriptsize\bfseries,text=prgray}
]
\node[stage,fill=prlightgold] (generate) at (0,0)
  {random process\\or controlled motif};
\node[tag,above=1mm of generate] {generate};

\node[stage,fill=prlightblue] (models) at (4.4,0)
  {two process-model\\realizations\\$N^{(1)},N^{(2)}$};
\node[tag,above=1mm of models] {construct};

\node[stage,fill=prlightteal] (observe) at (8.8,0)
  {relation evidence\\and two event log observations\\$L^{(u)},L^{(r)}$};
\node[tag,above=1mm of observe] {relate and observe};

\node[stage,fill=prlightred,minimum width=38mm] (rows) at (13.5,0)
  {canonicalize labels\\pair models with event logs\\\textbf{one family $\rightarrow$ four rows}};
\node[tag,above=1mm of rows] {align};

\draw[flow] (generate)--(models);
\draw[flow] (models)--(observe);
\draw[flow] (observe)--(rows);
\end{tikzpicture}%
}
\caption{Behavior-family construction in four stages. Family membership is established before row expansion, so both model realizations and both event log observations remain in the same data split.}
\label{fig:corpus}
\end{figure}
